\section{Reflexion und Grenzen}

\subsection{Erfüllung der Modulanforderungen (Datenbanken 2)}
\begin{itemize}
  \item \textbf{Öffentliche API als Datenquelle:} FMP API ✓ % Repo-Evidenz prüfen
  \item \textbf{Pandas-Verarbeitung:} Normalisierung, Gates, Akkumulation ✓ % Repo-Evidenz prüfen
  \item \textbf{MySQL als Datenbank:} Clean Store mit strukturierten Tabellen ✓ % Repo-Evidenz prüfen
  \item \textbf{MongoDB als Datenbank:} Raw Store für flexible Daten ✓ % Repo-Evidenz prüfen
  \item \textbf{Interaktive Streamlit-UI:} Dashboard, Explorer, Methodik ✓ % Repo-Evidenz prüfen
  \item \textbf{Mindestens eine Visualisierung:} Buy/Sell-Charts, Trends, KPIs ✓ % Repo-Evidenz prüfen
  \item \textbf{Beide Datenbanken gelesen:} MySQL für Analyse, MongoDB für Raw-Counts in Admin ✓ % Repo-Evidenz prüfen
  \item \textbf{Akademische Nachvollziehbarkeit:} Methodologie-Seite, Dokumentation ✓
\end{itemize}

\subsection{Technische Grenzen}
\begin{itemize}
  \item \textbf{API-Rate-Limits:} FMP hat tägliche Obergrenzen (üblicherweise 250 Calls/Tag). % Externe Quelle nötig
    \begin{itemize}
      \item Auswirkung: Nur eine begrenzte Anzahl von Trades und Profilen pro Tag möglich.
      \item Mitigation: Lokale Gates reduzieren API2/API3-Aufrufe um ~80\%. % Interne Projektspezifikation
    \end{itemize}
  \item \textbf{Datenqualität und Payload-Änderungen:} FMP kann API-Schemas ändern. % Externe Quelle nötig
    \begin{itemize}
      \item Raw Store ermöglicht Debugging und Reimport bei Änderungen.
      \item Normalisierung ist zentral; Feldmutationen erfordern Code-Änderungen.
    \end{itemize}
  \item \textbf{Identitätsprobleme (Ticker vs. CIK):} Mapping-Fehler zwischen Symbol und CIK sind möglich. % Repo-Evidenz prüfen
    \begin{itemize}
      \item Fallback-Logik (CIK-Suche) ist implementiert, aber nicht immer zuverlässig.
      \item Manuelle Korrektur in Admin möglich.
    \end{itemize}
  \item \textbf{Fehlende oder veraltete Profile:} Nicht alle Symbole haben aktuelle FMP-Profile. % Externe Quelle nötig
    \begin{itemize}
      \item Auswirkung: Fehlende market\_cap, industry reduzieren Score-Qualität.
      \item Workaround: Nachträgliches Laden über Admin-Button.
    \end{itemize}
  \item \textbf{Skalierbarkeit:} Single-Instance Streamlit, keine horizontale Skalierung ohne Docker/Cloud. % Eigene methodische Ableitung
    \begin{itemize}
      \item Für akademische MVP ausreichend.
      \item Für produktive Nutzung: mehrere Worker, Message Queue, caching erforderlich.
    \end{itemize}
\end{itemize}

\subsection{Methodische Grenzen}
\begin{itemize}
  \item \textbf{Signaltrennung nicht validiert:} % Eigene methodische Ableitung
    \begin{itemize}
      \item Score wurde nicht gegen historische Kursentwicklung validiert.
      \item Keine statistische Signifikanz-Tests.
      \item Akademische Baseline-Studien zu Insider-Trades nicht systematisch berücksichtigt. % Externe Quelle nötig
    \end{itemize}
  \item \textbf{Insider-Motive unklar:} % Externe Quelle nötig
    \begin{itemize}
      \item Insider handeln aus vielen Gründen (Liquidität, Steuer, Diversifizierung), nicht nur aus Information.
      \item Score kann diese Nutzenmotive nicht unterscheiden.
      \item Akademische Forschung zeigt: nur unter bestimmten Bedingungen sind Insider-Käufe prädiktiv. % Externe Quelle nötig
    \end{itemize}
  \item \textbf{Marktumgebung nicht berücksichtigt:} % Externe Quelle nötig
    \begin{itemize}
      \item Insider-Signale sind in Bull-Märkten und Bear-Märkten unterschiedlich stark.
      \item Das Score-Modell berücksichtigt die Marktphase nicht.
    \end{itemize}
  \item \textbf{Hintertests und Parameter-Justierung:} % Eigene methodische Ableitung
    \begin{itemize}
      \item Gewichtungen wurden nicht durch Backtesting kalibriert.
      \item Aktuell nur Cloudflare Quick Tunnel aktiv; Architektur ist auf weitere Provider (z. B. ngrok, Managed Tunnel) vorbereitbar.
    \end{itemize}
\end{itemize}

\subsection{Datenschutz und regulatorische Aspekte}
\begin{itemize}
  \item \textbf{Öffentliche SEC-Daten:} Alle Inputs sind SEC Form 4 Meldungen (öffentlich verfügbar). % Externe Quelle nötig
  \item \textbf{Keine Anlageberatung:} Mercator ist kein Finanzdienstleistungsprodukt. % Eigene methodische Ableitung
  \item \textbf{Disclaimer erforderlich:} Nutzer müssen verstehen, dass Score keine Prognose ist. % Eigene methodische Ableitung
\end{itemize}

\subsection{Interne Parameter -- kritische Würdigung}
\begin{itemize}
  \item \textbf{3-Tage-Akkumulations-Fenster:} % Interne Projektspezifikation
    \begin{itemize}
      \item Keine externe Begründung; gewählt für Projektkontext.
      \item Alternativen: 5, 7, 10 Tage sind ebenfalls vertretbar. % Eigene methodische Ableitung
      \item Impact: Größeres Fenster = stärkere Signale, aber breiter/weniger spezifisch.
    \end{itemize}
  \item \textbf{500-Tage-Lookback (API3):} % Interne Projektspezifikation
    \begin{itemize}
      \item Typisch für technische Analyse (1-2 Jahre Daten).
      \item Größere Fenster liefern bessere Trend-Estimates, erfordern aber mehr API-Aufrufe.
    \end{itemize}
  \item \textbf{Minimum Trade Value (100.000 USD):} % Interne Projektspezifikation
    \begin{itemize}
      \item Filtert Rauschen, kann aber echte Signale ausschließen.
      \item Konfigurierbar im Admin. % Repo-Evidenz prüfen
    \end{itemize}
  \item \textbf{Score-Gewichtungen:} % Interne Projektspezifikation
    \begin{itemize}
      \item Transaction Code 40\%, Trade Value 20\%, Market Cap 15\%, Trend 15\%, Age/Liquidity je 5\% (Beispiel).
      \item Nicht empirisch validiert; Adjustments möglich.
    \end{itemize}
\end{itemize}

\subsection{Verbesserungspotenziale}
\begin{itemize}
  \item \textbf{Robustere Retry- und Fehlerbehandlung:}
    \begin{itemize}
      \item Circuit-Breaker für API-Fehler.
      \item Detaillierte Error-Logging und Monitoring.
    \end{itemize}
  \item \textbf{Erweiterte Tests:}
    \begin{itemize}
      \item Unit-Tests für Gate-Logik und Scoring (teilweise vorhanden).
      \item End-to-End-Tests mit echten API-Daten (Playwright-Tests existieren). % Repo-Evidenz prüfen
      \item Integrationstests für Raw/Clean-Sync.
    \end{itemize}
  \item \textbf{Bessere Datenqualitäts-Berichte:}
    \begin{itemize}
      \item Dashboard für fehlende Profile, ungültige Trades, API-Fehler.
      \item Anomaly Detection (z. B. plötzliche Spitzen in Trade Volume).
    \end{itemize}
  \item \textbf{Erweiterte Visualisierungen:}
    \begin{itemize}
      \item Netzwerk-Graph (Insider $\leftrightarrow$ Company).
      \item Heatmaps für Sektor/Period-Kombinationen.
      \item Detaillierte Backtests und Performance-Analysen.
    \end{itemize}
  \item \textbf{Parallelisierung und Caching:}
    \begin{itemize}
      \item Asynchrone API-Aufrufe statt sequenziell.
      \item Redis-Cache für Profile und technische Metriken.
    \end{itemize}
  \item \textbf{Mobile Responsiveness:}
    \begin{itemize}
      \item Aktuelle Streamlit-Version hat gutes Mobile-Support; weitere Optimierungen möglich.
    \end{itemize}
\end{itemize}
